\documentclass[a4paper,14pt]{extarticle}
\usepackage[utf8x]{inputenc}
\usepackage[T1,T2A]{fontenc}
\usepackage[russian]{babel}
\usepackage{hyperref}
\usepackage{indentfirst}
\usepackage{listings}
\usepackage{color}
\usepackage{here}
\usepackage{array}
\usepackage{multirow}
\usepackage{graphicx}

%\usepackage{mathptmx}

%% Переименование "содержания" в "оглавление"
\renewcommand\contentsname{Оглавление}
\renewcommand\refname{Список использованных источников}

\bibliographystyle{ugost2008ls}

\usepackage{caption}
\renewcommand{\lstlistingname}{Программа} % заголовок листингов кода

\usepackage{listings}
\lstset{ %
extendedchars=\true,
keepspaces=true,
language=bash,					% choose the language of the code
basicstyle=\footnotesize,		% the size of the fonts that are used for the code
numbers=left,					% where to put the line-numbers
numberstyle=\footnotesize,		% the size of the fonts that are used for the line-numbers
stepnumber=1,					% the step between two line-numbers. If it is 1 each line will be numbered
numbersep=5pt,					% how far the line-numbers are from the code
backgroundcolor=\color{white},	% choose the background color. You must add \usepackage{color}
showspaces=false				% show spaces adding particular underscores
showstringspaces=false,			% underline spaces within strings
showtabs=false,					% show tabs within strings adding particular underscores
frame=single,           		% adds a frame around the code
tabsize=2,						% sets default tabsize to 2 spaces
captionpos=b,					% sets the caption-position to bottom
breaklines=true,				% sets automatic line breaking
breakatwhitespace=false,		% sets if automatic breaks should only happen at whitespace
escapeinside={\%*}{*)},			% if you want to add a comment within your code
postbreak=\raisebox{0ex}[0ex][0ex]{\ensuremath{\color{red}\hookrightarrow\space}}
}

%% Полуторный интервал
\usepackage[nodisplayskipstretch]{setspace}
\onehalfspacing

\usepackage[a4paper,includefoot,includehead,left=2.5cm,right=1.8cm,
top=2cm,bottom=2.5cm,bindingoffset=0cm]{geometry}


%% Выравнивание номара страницы по правому краю
\usepackage{fancyhdr}
\pagestyle{fancy}
\fancyhf{}
\protect\fancyfoot[R]{\thepage}
\renewcommand{\headrulewidth}{0pt}
\renewcommand{\footrulewidth}{0pt}

%% Нумерация картинок по секциям
\usepackage{chngcntr}
\counterwithin{figure}{section}
\counterwithin{table}{section}

%% Поля подписи и даты
\newcommand{\sign}[1][5cm]{%
\makebox[#1]{\hrulefill}
}

%% Количесво рисунков
\usepackage{totcount}
\newtotcounter{citenum}
\def\oldcite{}
\let\oldcite=\bibcite
\def\bibcite{\stepcounter{citenum}\oldcite}

\usepackage[figure,table,lstlisting]{totalcount}
% страниц
\usepackage{lastpage}
% библиографий
\newtotcounter{citnum} %From the package documentation
\def\oldbibitem{} \let\oldbibitem=\bibitem
\def\bibitem{\stepcounter{citnum}\oldbibitem}
% листингов (подключённых)
\newtotcounter{listnum}
\def\oldlstinputlisting{} \let\oldlstinputlisting=\lstinputlisting
\def\lstinputlisting{\stepcounter{listnum}\oldlstinputlisting}

%% Добавление промежуточных точек в оглавление
\usepackage{tocloft}
\renewcommand{\cftsecdotsep}{\cftdotsep}
\renewcommand{\cftsecleader}{\cftdotfill{\cftsecdotsep}}

%%Точки нумерации заголовков
\usepackage{titlesec}
\titlelabel{\thetitle.\quad}
\usepackage[dotinlabels]{titletoc}

%% Оформления подписи рисунка
\addto\captionsrussian{\renewcommand{\figurename}{Рисунок}}
\captionsetup[figure]{labelsep = period}

%% Подпись таблицы
\DeclareCaptionFormat{hfillstart}{\hfill#1#2#3\par}
\captionsetup[table]{format=hfillstart,labelsep=newline,justification=centering,skip=-10pt,textfont=bf}

%% Путь к каталогу с рисунками
\graphicspath{{figs/}}


\begin{document}	% начало документа
\counterwithin{lstlisting}{section}

% Титульная страница
%\begin{titlepage}	% начало титульной страницы

	\begin{center}		% выравнивание по центру

		\large Санкт-Петербургский Политехнический Университет Петра Великого\\
		\large Институт компьютерных наук и технологий \\
		\large Кафедра компьютерных систем и программных технологий\\[6cm]
		% название института, затем отступ 6см
		
		\huge \textbf{КУРСОВОЙ ПРОЕКТ}\\[0.5cm]
		\large Название предмета\\[0.1cm]
		\large Тема работы\\[5cm]

	\end{center}


	\begin{flushright} % выравнивание по правому краю
%		\begin{minipage}{0.5\textwidth} % врезка в половину ширины текста
%			\begin{flushleft} % выровнять её содержимое по левому краю

				\large Выполнил студент группы 43501/4\\
				\large В.Д. Петров\\[0.5cm]
				
				\large Принял к.т.н., доцент\\
				\sign[4cm]\large  В.М. Ицыксон\\
				\large Оценка: \sign\\
				«\underline{\hspace{0.7cm}}» \underline{\hspace{2cm}} \the\year г.

%			\end{flushleft}
%		\end{minipage}
	\end{flushright}
	
	\vfill % заполнить всё доступное ниже пространство

	\begin{center}
	\large Санкт-Петербург\\
	\large \the\year % вывести дату
	\end{center} % закончить выравнивание по центру

\thispagestyle{empty} % не нумеровать страницу
%\end{titlepage} % конец титульной страницы
\newpage


% Реферат

\keywords{%
  ТРАНСФОРМАЦИЯ ПРОГРАММ,
  СТАТИЧЕСКИЙ АНАЛИЗ,
  КАЧЕСТВО ПРОГРАММНОГО ОБЕСПЕЧЕНИЯ
}

\abstractcontent{
Краткое описание работы
}


% Содержание
% Содержание
\renewcommand\contentsname{\centerline{Содержание}}
\tableofcontents
\thispagestyle{fancy}
\newpage



\section*{Введение}
\addcontentsline{toc}{section}{Введение}


\section{Программа работы}


\section{Теоретическая информация}


\section{Ход выполнения работы}

\subsection{Список}

\begin{itemize}
\item первый элемент списка
\item второй элемент списка
\end{itemize}


\subsection{Картинка}

\begin{figure}[H]
	\begin{center}
		\includegraphics[scale=0.7]{sample}
		\caption{название картинки} 
		\label{pic:pic_name} % название для ссылок внутри кода
	\end{center}
\end{figure}


\subsection{Листинг}

\captionof{lstlisting}{hell\_o.c} % для печати символ '_' требует выходной символ '\'
\lstinputlisting[label=code:hello]{listings/hell_o.c}
\parindent=1cm % командна \lstinputlisting сбивает параментры отступа
Текст без отступа (следует за вставкой)

Новый параграф

\noindent Новый параграф с принудительно выключенным отступом


\subsection{Частичный листинг}
% настрока частичного ввода (требуется один раз)
\makeatletter
\def\lst@PlaceNumber{\llap{\normalfont
                \lst@numberstyle{\the\lst@lineno}\kern\lst@numbersep}}
\makeatother

\captionof{lstlisting}{фрагмент hell\_o.c}
\lstinputlisting[label=code:hello_mod, linerange={4-5}]{listings/hell_o.c}
\parindent=1cm

\subsection{Таблица}

\begin{table}[H]
	\caption{Название таблицы}
	\begin{center}
		\begin{tabular}{|l|l|}
			\hline
			top left & top right\\ \hline
			bot left & bot right\\ \hline
	\end{tabular}
		\label{tabular:tab_examp}
	\end{center}
\end{table}

\section*{Заключение}
\addcontentsline{toc}{section}{Заключение}
\LaTeX\ удобен для создания отчётов, так как сам следит за нумерацией таблиц, рисунков, листингов и отсылок к ним (так, например, здесь всегда будет указан номер рисунка "sample" не зависимо от того, какой он (1,2 или другой) - это рисунок \ref{pic:pic_name}). Не менее важно что весь документ оформлен в едином стиле, а исходные материалы подключаются к отчёту, а не хранятся в нём. Всё это позволяет легко получить качественный отчёт без дополнительных трат на его офрмление.

Исключения, пожалуй, составляют таблицы, так как их значительно сложнее создавать кодом, нежели в графическом редакторе. Но здесь никто не запрещает использовать визуальные средства создания таблиц для \LaTeX\ . А пока всё само такое умное, можно почитать хорошую литературу \cite{saturday_is_monday}.


\renewcommand\refname{Список использованных источников}
\bibliography{sources}
\addcontentsline{toc}{section}{Список использованных источников}

\end{document}
